\documentclass[11pt]{article}
\usepackage[a4paper,margin=1in]{geometry}
\usepackage{amsmath}

\title{folio sample document}
\author{}
\date{}

\begin{document}
\maketitle

This is a multi-page sample used to exercise folio's rendering and page
navigation. The same content is provided in three parallel formats --- Markdown,
LaTeX, and PDF --- so you can compare how each input path looks. Page with
\texttt{n} / \texttt{p}, jump with \texttt{:n}, and zoom with \texttt{+} /
\texttt{-}.

\section{Introduction}
folio renders document pages as images directly in your terminal using the sixel
or kitty graphics protocols. It opens PDFs natively, converts Markdown to PDF on
the fly, and---when you opt in via the \texttt{latex\_engine} setting---compiles
LaTeX with an installed TeX engine. Whatever the source, the same viewer provides
zoom, pan, dark mode, configurable sharpness, bookmarks, and resume.

This document was typeset by a real TeX engine and then rasterized by folio. The
text is intentionally long enough to span more than two pages so that multi-page
behaviour is easy to try out.

\section{Text formatting}
Paragraphs may contain \textbf{bold}, \textit{italic}, \textbf{\textit{bold
italic}}, and inline \texttt{code} such as \texttt{let x = 42;}. Long paragraphs
wrap to the page width and are justified by the TeX engine, breaking cleanly and
hyphenating where appropriate rather than running past the margin.

Source line breaks are collapsed into the same paragraph, so you can keep your
\LaTeX{} lines short without affecting the rendered output; a blank line starts a
new paragraph.

\section{Lists}
Unordered lists, including nested items:
\begin{itemize}
    \item The quick brown fox jumps over the lazy dog.
    \item A second item with a little more text so it wraps to demonstrate
          hanging indentation under the bullet.
    \begin{itemize}
        \item A nested item.
        \item Another nested item.
    \end{itemize}
    \item A third top-level item.
\end{itemize}

Ordered lists:
\begin{enumerate}
    \item Preheat the oven.
    \item Mix the ingredients thoroughly.
    \item Bake until golden, then let cool before serving.
\end{enumerate}

\section{Code}
Verbatim text renders in a monospace face:
\begin{verbatim}
fn main() {
    let pages = vec!["one", "two", "three"];
    for (i, name) in pages.iter().enumerate() {
        println!("page {}: {}", i + 1, name);
    }
}
\end{verbatim}

\section{A longer discussion}
Terminal document viewers occupy a useful niche: they keep you in the
keyboard-driven, text-first environment where so much development happens, while
still showing real, faithfully rendered pages. Rather than dumping a lossy text
approximation of a document, folio rasterizes the genuine page and paints it with
a graphics protocol, so figures, fonts, and layout survive intact.

Because only the page under view is rendered on demand, memory and CPU stay
modest even for large documents, and an LRU cache plus predictive prefetch keep
navigation feeling immediate. The same machinery that makes PDFs responsive
applies to converted Markdown and LaTeX, since every input is normalised to PDF
before it reaches the renderer.

This uniformity is the point of the modular input layer. Adding a new input
format is a matter of teaching folio how to turn that format into PDF bytes; the
viewer, the cache, the themes, and the bookmark store do not need to know where
those bytes came from. Markdown was the first such format, LaTeX the second, and
the same seam is ready for whatever comes next.

\section{Mathematics}
Unlike the Markdown path, the LaTeX path typesets mathematics properly. Inline
math such as $E = mc^2$ and $\sum_{i=1}^{n} i = \frac{n(n+1)}{2}$ sits naturally
within the prose, with Greek letters like $\alpha$, $\beta$, and $\pi$ rendered as
real glyphs.

A displayed equation:
\begin{equation}
    \int_0^1 x^2 \, dx = \frac{1}{3},
    \qquad
    e^{i\pi} + 1 = 0.
\end{equation}

\section{Conclusion}
If you can see this section, pagination is working and you have paged past the
first page. Try toggling dark mode with \texttt{i}, adjusting sharpness with
\texttt{[} and \texttt{]}, dropping a bookmark with \texttt{m}, and reopening the
file to confirm it resumes where you left off.

\end{document}
